{\large\textbf{E2 - Optical black-box (10 pts)}}

{\LARGE WARNING!}

\noindent\fbox{\parbox{\dimexpr\linewidth-2\fboxsep-2\fboxrule\relax}{\textbf{Do
not open the black-box. Do not shake the black-box. Do not touch the windows of
the optical ports. If you break the black-box or the windows or attempt to open
the black box, you will be disqualified.}}}

Your task is to determine the contents of an optical black box without opening
it.

The black-box has four optical ports (A, B, C and D) for light, and two optical
axes (Fig. 2). The optical axes are perpendicular to each other. There is up to
one optical element behind each of the ports as well as another one in the
center of the box. You can use a laser and a laser mount with a wheel (that you
can put marks on it with a pencil) in order to rotate the laser.

\textbf{Equipment (see also Fig. 4)}

A\quad Black-box

B\quad Laser module in mount with wheel (the same laser module is used for both
experiments, to be placed on the table surface)

C\quad Transparent block

D\quad Masking tape, pencil, ruler, paper with diagonal scale

\noindent\fbox{\parbox{\dimexpr\linewidth-2\fboxsep-2\fboxrule\relax}{\textbf{Do
not look directly into the laser beam, and make sure no other people are hit by
it. Do not look into the optical ports of the box if the laser is on, and turn
the laser off when not needed.}}}

\textbf{Task E2.1 - Central element ($\sim$0.3 pts)}

The two optical axes cross in the center of the black box. At the crossing
could sit: no element, a fully reflective mirror (both sides), a
semi-transparent mirror, or a regular-triangle-shaped prism.

Table 1: Possible elements in the slots of the black-box

\begin{center}
\begin{tabular}{>{\raggedright\arraybackslash}p{2.6cm}p{6cm}}
\hline
no element & there is just air in the slot \\
\hline
mirror & angle between the mirror axis and one of the optical axes \\
\hline
prism & angle between one of its sides and one of the optical axes of the black
box, shaped like regular triangle \\
\hline
concave or convex lens & distance to the center of the box, magnitude and sign
of the focal length. \textit{Note: Axes of lenses are always along optical
axes.} \\
\hline
polarizer & angle of orientation relative to the vertical axis of the black-box \\
\hline
single, thin slit & distance to the center of the box, width of the slit \\
\hline
diffraction grating & distance to the center of the box, direction of stripes,
distance between the stripes \\
\hline
pin hole & distance to the center of the box, hole diameter \\
\hline
\end{tabular}
\end{center}

Determine which element is placed centrally in the black box. Describe its
orientation towards the optical ports (A, B, C, D) - for example by using a
sketch. Justify your choices.

\begin{center}\includegraphics[width=0.42\linewidth]{figures/EuPhO_2023_Q5_fig1.png}\end{center}

Figure 2: Layout of the black box and the slots of the unknown elements

\textbf{Task E2.2 - Elements in remaining slots ($\sim$2.2 pts)}

There is one element from Table 1 in each of the four slots behind the optical
ports A, B, C, D, respectively.

Determine for each slot the type of element present. Justify your choices.

\textbf{Task E2.3 - Properties ($\sim$7.5 pts)}

In Table 1, you can find a second column containing characteristic properties
of the possible elements.

Determine these characteristic properties for the optical elements used inside
the box at slots A, B, C and D \textbf{as precisely as possible}.

\textbf{Important hints:}

\textbullet\ The wavelength of the laser is $(650 \pm 5)\,\mathrm{nm}$.

\textbullet\ The refractive index of transparent elements can be assumed to be
$1.5$.

\begin{center}\includegraphics[width=0.55\linewidth]{figures/EuPhO_2023_Q5_fig2.jpg}\end{center}

Figure 4: Setup and equipment for experimental problem E2. \textit{Note: You can
secure the laser module to the table using the tape - see B.}
